\documentclass[11pt]{article}
\usepackage{amsmath,amssymb,amsthm,mathtools}
\usepackage[margin=1in]{geometry}
\theoremstyle{plain}
\newtheorem{theorem}{Theorem}
\newtheorem{lemma}{Lemma}
\theoremstyle{remark}
\newtheorem{remark}{Remark}
\newcommand{\E}{\mathbb{E}}\newcommand{\R}{\mathbb{R}}

\title{Formalization brief: \texttt{lem\_gauss} as the single-field instance\\
of the proved martingale-problem architecture}
\date{8 July 2026}

\begin{document}
\maketitle

\noindent\textbf{Goal.} Remove the \texttt{sorry} at \texttt{lem\_gauss} in
\texttt{TypeDDecouplingEW.lean} (attached). Two observations drive the design:
(1) the paper itself derives Lemma 6.6 as ``the single-field case of
Theorem~6.2,'' and Theorem 6.2's Gaussian machinery is now PROVED
(\texttt{TypeDDecouplingMartingaleGaussian.lean}, attached; do not modify);
(2) the current statement is a FIDELITY DEFECT (document as catch \#8): it
takes a free $Y$ with no dynamical hypotheses and asserts OU convergence ---
a false-universal-shaped placeholder, not a rendering of Lemma 6.6. Whole
project must build; standard axioms only.

\section*{Part 0: fidelity repair (mandatory)}

Restate \texttt{lem\_gauss} with the faithful single-species hypotheses tying
$Y$ to the WASEP dynamics, mirroring \texttt{thm\_mp}'s de-opaqued
architecture: a single-field drift condition (the (D) analogue of
\texttt{mpConvDrift}), a single-field bracket-realization condition (the (N)
analogue of \texttt{mpConvBracket}: the pairings are charFun-realized as
martingale-difference arrays with deterministic bracket
$2\chi Dt\cdot\mathrm{sig}(\varphi)$, $\chi=\rho(1-\rho)$, with stopped
companions), and the path-space bundle (an \texttt{MPPathBundle}-style field).
Document: the previous free-$Y$ form did not assert the paper's lemma; the
Dittrich--G\"artner citation moves to the docstring as the classical theorem
this instantiates.

\section*{Part 1: prove the restated lemma}

From the restated hypotheses, the conclusion
$\exists Y_{\mathrm{lim}}$, \texttt{convInLawDist} $\wedge$
\texttt{isStationaryOU} follows by the SAME derivation as \texttt{thm\_mp}:
\texttt{mpConvBracket\_gaussian\_fdd}-style invocation of
\texttt{martingale\_charFun\_gaussian} (single field --- no joint/cross
component needed) plus the path bundle. Reuse the \texttt{thm\_mp} proof
structure; if a shared lemma serves both, factor it rather than duplicate.

\section*{Part 2: the WASEP (N)-computation --- the genuinely new content}

Prove, as standalone library-clean lemmas (suggest
\texttt{TypeDDecouplingBracketN.lean}, or extend the Drift file's namespace
style; do NOT modify existing files), the finite-$N$ facts that make the
bracket hypothesis faithful for the actual single-species WASEP:
\begin{itemize}
\item[(a)] \underline{Mean:} for the single-species exclusion with rates
$(1,q^2)$ under a Bernoulli($\rho$) product weight on the window, the
integrated instantaneous bracket of the field pairing --- the normalized sum
$\frac1N\sum_x\varphi'(x/N)^2\,\E[c_x(\eta)(\nabla_x\text{-increment})^2]$
with $c_x$ the exclusion jump rates --- equals the explicit Riemann sum
converging to $2\chi D\|\varphi'\|_{L^2}^2$ (compute the exact per-bond
equilibrium factor: $\E[c_x] = (1+q^2)\rho(1-\rho)$-type; identify
$D=(1+q^2)/2$-normalization consistently with \texttt{prop\_drift}'s $D$ ---
derive, don't guess; any consistent convention documented).
\item[(b)] \underline{Variance:} the fluctuation of the bracket functional has
second moment $O(1/N)$ under the product weight --- the
\texttt{corr\_second\_moment} pattern (attached Drift file): local functions,
finite-range covariances, $(1/N^2)\cdot O(N)$ bookkeeping.
\item[(c)] \underline{Time integration:} from (a)+(b), the time-integrated
bracket over $[0,t]$ concentrates at $2\chi Dt\|\varphi'\|^2$ in $L^2$ under
stationarity (equal-time Cauchy--Schwarz in time, as in \texttt{prop\_drift}'s
correction term --- no mixing needed).
\end{itemize}
These are [P]-class results consumed by the restated hypothesis's
documentation: the bundle field asserting the bracket realization cites (a)--(c)
as the finite-$N$ ground truth, with only the Dynkin-bracket identification
(the \texttt{dynkinBracketDef} convention) and the process-level pairing
construction remaining as the documented citation content.

\section*{Part 3: wiring}

\texttt{thm\_ewmain} consumes \texttt{lem\_gauss} for its diagonal-noise
input: thread the new hypotheses exactly as \texttt{hmp}/\texttt{hconc} are
threaded (add bundle parameters, pass through). No other leaf touched;
\texttt{thm\_mitoma}/\texttt{prop\_aldous} unchanged. Report: the restated
hypothesis list, what (a)--(c) proved, the bundle contents, and the final
count (target 3 $\to$ 2).

\begin{remark}[hygiene]
(1) The restated hypotheses must be SATISFIABLE and faithful --- each field a
standard fact of the stationary WASEP field; say so per field, citing the
paper location (Lemma 6.6's proof paragraph and \S6.1).
(2) Constants/conventions: match \texttt{prop\_drift}'s $D$ and the paper's
$2\chi D$ normalization; if the encoding's $\mathrm{sig}(\varphi)$ plays the
role of $\|\varphi'\|^2$, keep it as data and state (a) against it.
(3) Do not modify \texttt{TypeDDecouplingMartingaleGaussian.lean},
\texttt{TypeDDecouplingDrift.lean}, or the CLT files; new lemmas go in the new
file or the EW file.
\end{remark}

\end{document}
